% Paper 18 §IV drop-in: squared Dirac operator and the completeness of the taxonomic grid
% Insert after the "Certification and taxonomy wiring" paragraph, before "Parallel to Paper 24"
% Track T9, Tier 3 sprint, 2026-04-15

\paragraph{The squared Dirac operator and the completeness of the taxonomic grid.}
The Tier 3 Dirac-on-$S^3$ sprint (2026, Track T9) probed the last open
cell in the operator-order $\times$ bundle classification
(Table~\ref{tab:bundle_operator_tiers}): the squared Dirac operator
$D^2$ on the unit three-sphere, a second-order operator on spinor
sections.

By the Lichnerowicz formula~\cite{lichnerowicz1963}, $D^2 = \nabla^*\nabla + R/4$,
where $R = 6$ is the scalar curvature of unit $S^3$, so $R/4 = 3/2$ is a
rational constant shift.  The eigenvalues of $D^2$ are
$(n + 3/2)^2 = (2n+3)^2/4$ with degeneracies $g_n = 2(n+1)(n+2)$ (full
Dirac, both chiralities).  The spectral zeta function evaluates in closed
form at every integer $s$ via a partial-fraction identity: substituting
$m = n+1$ and using $m(m+1) = [(2m+1)^2 - 1]/4$,
\begin{equation}\label{eq:zeta_d2}
  \zeta_{D^2}(s)
  \;=\; 2^{2s-1}\bigl[\lambda(2s{-}2) - \lambda(2s)\bigr],
\end{equation}
where $\lambda(a) = (1 - 2^{-a})\,\zeta_R(a)$ is the Dirichlet lambda
function.  Since $\zeta_R(2k) = (-1)^{k+1}\,(2\pi)^{2k}\,B_{2k}/(2\cdot(2k)!)$
(Bernoulli formula), every term is a rational multiple of $\pi^{2k}$.
At every integer $s$, $\zeta_{D^2}(s)$ is therefore a two-term polynomial
in $\pi^2$ with rational coefficients:
$c_{2s-2}\,\pi^{2s-2} + c_{2s}\,\pi^{2s}$.  Explicitly:
$\zeta_{D^2}(1) = -\pi^2/4$ (analytic continuation),
$\zeta_{D^2}(2) = \pi^2 - \pi^4/12$,
$\zeta_{D^2}(3) = \pi^4/3 - \pi^6/30$,
$\zeta_{D^2}(4) = \pi^6(168 - 17\pi^2)/1260$.

\emph{No odd-zeta content appears at any~$s$.}  The structural mechanism
is algebraic: $D^2$ eigenvalues are perfect squares of odd half-integers,
so the spectral sums reduce to $\sum_k f(k)\,k^{-2s}$ over odd integers,
with exponent $2s$ (always even).  By the Bernoulli formula, such sums
are always $\mathbb{Q} \cdot \pi^{2s}$.  Squaring the operator maps the
first-order Dirichlet $\lambda(\text{odd})$ to $\lambda(\text{even})$,
structurally eliminating the odd-zeta content that Tier~1 (Track~D3)
found in the first-order Dirac spectral zeta via
Eq.~(\ref{eq:dirac_dirichlet_zeta3}).  The Hopf-equivariant per-charge
decomposition (half-integer charges for spinor modes) confirms this
universally across all Hopf sectors.

The (2nd-order, spinor-bundle) cell of Table~\ref{tab:bundle_operator_tiers}
is therefore \textbf{filled and degenerate with the scalar calibration
cell}: both produce $\pi^{\text{even}}$ content exclusively.  The
Lichnerowicz constant shift $R/4 = 3/2$ modifies rational coefficients
but does not introduce new transcendental content.  Operator order---first
vs.\ second---is the primary discriminant for transcendental class on
round~$S^3$.  Bundle type (scalar vs.\ spinor) modulates the
degeneracy lattice ($n^2$ vs.\ $2(n+1)(n+2)$), the eigenvalue lattice
(integer vs.\ half-integer), the Hopf charge lattice (integer vs.\
half-integer), and the rational coefficients in the $\pi^{2k}$
expansions, but it does \emph{not} change the transcendental class.

The taxonomic grid thus collapses to three effective tiers:
\begin{enumerate}
  \item \textbf{First-order operators:} produce odd-zeta content
    ($\zeta_R(3), \zeta_R(5), \ldots$) via sums over
    $(\text{odd})^{-s}$ at odd~$s$, plus coupling content ($\alpha^2$,
    $\gamma$) in the spinor sector.
  \item \textbf{Second-order operators:} produce $\pi^{\text{even}}$
    content exclusively, on both scalar and spinor bundles. This is the
    calibration tier of Paper~24.
  \item \textbf{Basis and composition tiers:} unchanged (embedding,
    flow, composition constants are operator-independent).
\end{enumerate}
